% =====================================================================
%  diagramas.tex
%  Los cuatro diagramas del proceso de negocio RedPepper (F-04).
%  Las imágenes se reproducen sin retoques: conservan los rótulos
%  originales del material entregado por el equipo.
% =====================================================================

\section{Diagramas del proceso de negocio}
\label{sec:diagramas}

Esta sección reproduce los cuatro diagramas que describen la operación de
RedPepper. Conforme al reporte de origen (F-04), su contenido es
descripción complementaria de la operación, el abastecimiento y el cierre;
no constituye una nueva entrevista ni una fuente adicional de requisitos.

La Figura~\ref{fig:venta-mostrador} muestra el flujo de atención en
mostrador, desde la llegada del cliente hasta la entrega del pedido,
separando las responsabilidades de cliente, cajero, sistema y cocina. La
Figura~\ref{fig:abastecimiento} describe el abastecimiento y el prorrateo:
la compra registrada se clasifica como semanal o adicional y sigue caminos
de costo distintos antes de almacenar el comprobante. La
Figura~\ref{fig:cierre-mensual} presenta el cierre mensual, desde la
consolidación del margen bruto hasta la generación del reporte bajo
estructura de operación, financiamiento y fiscal. Finalmente, la
Figura~\ref{fig:proceso-general} integra los tres flujos anteriores en un
diagrama general con carriles para cliente, cajero, sistema, cocina, Jerry
como administrador, proveedor y contador.

\paragraph{Advertencia de lectura.} Los diagramas son material de diseño
propuesto por el equipo, no evidencia de un sistema ya operativo. Varias de
sus casillas (por ejemplo, el umbral de 48 px o el reparto entre cinco días
hábiles) coinciden con propuestas que el reporte de ilación identifica como
pendientes de validación con Jerry y con el contador
(Sección~\ref{sec:especificacion}).

\begin{figure}[H]
  \centering
  \includegraphics[width=\linewidth]{01-venta-mostrador.png}
  \caption{Diagrama 1: venta en mostrador. Flujo de atención con carriles
  de cliente, cajero, sistema y cocina; la comanda se genera al registrar
  la venta y cocina prepara al instante, sin productos precocidos.}
  \label{fig:venta-mostrador}
\end{figure}

\begin{figure}[H]
  \centering
  \includegraphics[width=\linewidth]{02-abastecimiento-prorrateo.png}
  \caption{Diagrama 2: abastecimiento y prorrateo. Tras registrar la compra
  se decide su tipo: la compra adicional se carga al 100\,\% en la fecha de
  ejecución y la semanal se prorratea entre cinco días hábiles; en ambos
  casos se adjunta el comprobante y se almacena cifrado con SHA-256.}
  \label{fig:abastecimiento}
\end{figure}

\begin{figure}[H]
  \centering
  \includegraphics[width=\linewidth]{03-cierre-mensual-niif18.png}
  \caption{Diagrama 3: cierre mensual NIIF 18. Del margen bruto mensual se
  restan los costos indirectos fijos para obtener el BAI; con las reglas
  tributarias que aporta el contador se calculan impuestos y cargas
  financieras, y se genera el reporte NIIF 18 para revisión remota.}
  \label{fig:cierre-mensual}
\end{figure}

\begin{figure}[H]
  \centering
  \includegraphics[width=\linewidth]{04-proceso-general.png}
  \caption{Diagrama general del proceso de negocio de RedPepper. Integra
  venta en mostrador, preparación en cocina, abastecimiento con proveedor,
  prorrateo de compras, validación tributaria del contador y cierre
  mensual con revisión remota del administrador.}
  \label{fig:proceso-general}
\end{figure}
